\chapter{RG, trace anomaly и абсолютный scale-setting}

Base-\(K\) renormalization gate оставил два продолжения: принять один
scale input либо вывести finite subtraction из RG fixed point или anomaly.
Проверим второй путь на поле-содержании уже построенного relative-\(U(1)\)
parent block.

\section{One-loop gauge beta coefficient}

Физический charge-square trace finite fermion module равен
\begin{equation}
 \sum_{\mathrm{Weyl}}q_f^2
 =\Tr_{\mathrm{orb}}Q^2=2.
\end{equation}
Два комплексных поля \(x,z\) имеют заряды \(-1,+1\), поэтому
\begin{equation}
 \sum_{\mathrm{complex\ scalar}}q_s^2=2.
\end{equation}
В стандартной четырёхмерной абелевой нормировке
\begin{equation}
 b
 =
 \frac23\sum_{\mathrm{Weyl}}q_f^2
 +\frac13\sum_{\mathrm{complex\ scalar}}q_s^2
 =\frac43+\frac23=2.
\end{equation}
Следовательно,
\begin{equation}
 \beta_g
 =\mu\frac{dg}{d\mu}
 =\frac{2}{16\pi^2}g^3
 +O(g^5).
\end{equation}

\begin{proposition}[One-loop fixed-point no-go]
В минимальном relative-\(U(1)\) field content единственная perturbative
one-loop fixed point равна \(g=0\). Значение \(g^2=3/8\) не является
fixed point; оно является spectral matching condition на некотором scale
\(\mu_{\mathrm{spec}}\).
\end{proposition}

\section{Running from the spectral matching condition}

One-loop решение имеет вид
\begin{equation}
 \frac1{g^2(\mu)}
 =
 \frac1{g^2(\mu_{\mathrm{spec}})}
 -\frac{b}{8\pi^2}
 \log\frac{\mu}{\mu_{\mathrm{spec}}}.
\end{equation}
При
\begin{equation}
 g^2(\mu_{\mathrm{spec}})=\frac38,
 \qquad b=2,
\end{equation}
формальная Landau scale удовлетворяет
\begin{equation}
 \log\frac{\Lambda_L}{\mu_{\mathrm{spec}}}
 =
 \frac{8\pi^2}{b\,g^2(\mu_{\mathrm{spec}})}
 =\frac{32\pi^2}{3}.
\end{equation}
Таким образом, spectral normalization фиксирует большое безразмерное
отношение
\begin{equation}
 \frac{\Lambda_L}{\mu_{\mathrm{spec}}}
 =\exp\left(\frac{32\pi^2}{3}\right),
\end{equation}
но не определяет абсолютное значение \(\mu_{\mathrm{spec}}\).

\section{Общая transmutation-константа}

Для автономного RG-уравнения
\begin{equation}
 \mu\frac{dg}{d\mu}=\beta(g)
\end{equation}
RG-invariant scale задаётся
\begin{equation}
 \Lambda_{\mathrm{RG}}
 =
 \mu
 \exp\left(
 -\int^g\frac{dg'}{\beta(g')}
 \right).
\end{equation}
Это выражение не вычисляет абсолютный scale из dimensionless
\(\beta\)-функции. Оно является интеграционной константой RG-trajectory,
определяемой одним boundary condition.

Сдвиг
\begin{equation}
 \mu\longmapsto s\mu,
 \qquad
 \Lambda_{\mathrm{RG}}\longmapsto s\Lambda_{\mathrm{RG}}
\end{equation}
оставляет все dimensionless running relations неизменными. Поэтому
dimensional transmutation заменяет безразмерный coupling на один
размерный вход, но не устраняет необходимость этого входа.

\begin{theorem}[RG integration-constant no-go]
Автономные beta-функции текущего scale-free parent action не могут
зафиксировать абсолютный \(\chi_*\). Они могут предсказывать отношения
scale-ов после задания одного matching scale, но сам matching scale
остаётся RG integration constant.
\end{theorem}

\section{Trace anomaly}

Trace anomaly имеет структуру
\begin{equation}
 T^\mu{}_\mu
 =
 \sum_i\beta_i(g)\,\mathcal O_i
 +\text{curvature anomalies}.
\end{equation}
Она описывает нарушение classical scale invariance и порождает
\(\Lambda_{\mathrm{RG}}\), но не задаёт boundary value RG-trajectory.
Топология \(K\), spin-структуры и charge lattice фиксируют operator
content, однако не выбирают начало logarithmic scale без дополнительного
normalization condition.

\begin{nogocriterion}[Anomaly-scale no-go]
Идентификация anomaly-generated \(\Lambda_{\mathrm{RG}}\) с физическим
mass scale не считается предсказанием, если не выведено независимое
условие, фиксирующее \(\Lambda_{\mathrm{RG}}R\) или эквивалентное
dimensionless отношение.
\end{nogocriterion}

\section{Вердикт}

\begin{theorem}[Scale-input inevitability in Version III]
В текущем minimal field content ни one-loop fixed point, ни dimensional
transmutation, ни trace anomaly не устраняют один абсолютный scale-setting
datum. Значение \(g^2=3/8\) остаётся строгим spectral matching relation,
а не scale-independent константой.
\end{theorem}

Следовательно, минимальная честная версия III должна:
\begin{enumerate}
 \item принять один mass или length scale как train input;
 \item вывести из parent action все dimensionless ratios и relative
 normalizations;
 \item зарезервировать минимум два независимых dimensionless observables
 для blind two-sector gate.
\end{enumerate}

\begin{protocol}[Следующий этап]
Необходимо выбрать один допустимый scale-setting input без оптимизации
по blind observables и построить таблицу:
\begin{equation}
 \text{один train scale}
 \longrightarrow
 \text{dimensionless predictions}
 \longrightarrow
 \text{два независимых blind tests}.
\end{equation}
Это переводит программу от невозможного absolute-scale no-input требования
к стандарту предсказательной effective theory с одной размерной
калибровкой.
\end{protocol}

Машинный аудит сохранён в
\path{s2t_v3_rg_anomaly_scale_setting_results.json}.